% ============================================================================
%  Chapter 6 -- Discussion and Conclusion
%  Evidence hierarchy: (A) main benchmark; (B) controlled ablation/refinement;
%  (C) mechanistic diagnostics; (D) calibration diagnostics; (E)
%  historical/supporting runs. A lower tier is never allowed to override or
%  silently reinterpret a higher tier's result.
% ============================================================================

\chapter{Discussion and Conclusion}
\label{ch:discussion-conclusion}

\section{Main Findings}
\label{sec:ch6-main-findings}

Chapter 5's primary benchmark (CIFAR-100 $5\times20$, seed 42)
places SimpleAvg's best result (SimpleAvg + KD, 75.21\%) 4.45 percentage
points above RankExt's best result (RankExt + FactorOrth + KD + Protect,
70.76\%). Two observations drive the rest of this chapter. First,
SimpleAvg's four variants move very little around a strong plain baseline
(75.16\%): none of KD, FactorOrth, or their combination produces a clear
all-seen improvement. Second, RankExt's four variants move an order of
magnitude more (34.99\% to 70.76\%), and the movement is concentrated almost
entirely in the presence or absence of the KD+Protect configuration
($+$30.02 percentage points over plain), with FactorOrth contributing a
smaller, separate increment ($+$6.01 percentage points alone; a further
$+$5.75 points on top of KD+Protect). This chapter asks why these two
families respond so differently to auxiliary objectives that are, at the
level of Chapter~\ref{ch:methodology}'s design, addressing the same nominal
problem --- retention under sequential class-incremental training --- and
what the parameter-efficiency, $R_{i,j}$, and refinement evidence of
Chapter~\ref{ch:results} imply about that difference. Every interpretation
below is explicitly graded by the evidence tier it rests on: (A) the
primary benchmark; (B) the controlled refinement experiment;
(C) mechanistic diagnostics; (D) calibration diagnostics; (E)
historical/supporting runs. A lower tier is never allowed to override or
silently reinterpret a higher tier's result.

\section{Why SimpleAvg Is Strong}
\label{sec:ch6-simpleavg-strong}

SimpleAvg's plain baseline reaches 75.16\% without any auxiliary mechanism
(Chapter~\ref{ch:results}, Section~\ref{sec:ch5-main-results}). Several
structural properties of its design plausibly contribute to this, and this
section states them as candidate contributors, not as an isolated, proven
cause, following a dedicated evidence audit specifically constructed to
test for leakage or an unfair evaluation advantage.

That audit confirms directly from the training and evaluation code --- not
inferred --- that SimpleAvg's final evaluation is genuinely task-agnostic: a
single merged model with a single 100-way classifier head, evaluated with a
plain argmax and no task-ID routing at inference, identical in mechanism to
RankExt's own evaluation path. No leakage of task identity into final
inference was found. Given that baseline of legitimacy, several structural
properties are consistent with SimpleAvg's strength: each of its five
specialists is trained independently from a fresh pretrained checkpoint, so
no task's gradient can interfere with another task's already-learned
parameters within a single evolving model --- the sequential interference
RankExt's architecture is exposed to by construction does not arise for
SimpleAvg at all. Classifier-row stitching (Chapter~\ref{ch:methodology})
assigns each class's final row directly from the one specialist that
actually trained it, preserving an exact, task-co-adapted linear boundary
rather than blending or averaging classifier weights across tasks. This
merge-then-stitch strategy sits within the broader model-merging literature
on combining independently fine-tuned weights
\parencite{wortsman2022model,ilharco2023editing}. Fresh, full-rank-80 LoRA
specialists and a $\times10$ classifier-head learning-rate multiplier
(unchanged throughout the primary benchmark) are both consistent with,
though not individually isolated as necessary for, SimpleAvg's result.

A specific, measured capacity question is worth stating precisely rather
than left implicit. An SVD-based effective-rank measurement on real,
architecture-matched checkpoints (an earlier diagnostic run for SimpleAvg, the calibration-diagnostic checkpoint for
RankExt --- the primary benchmark's own checkpoints were unavailable locally for this
measurement, a disclosed proxy) found that SimpleAvg's merged dense delta
has 137--192 singular values exceeding 1\% of the maximum, versus 79--80
for a single specialist's own delta (roughly $1.7\times$--$2.4\times$), and
an entropy-based effective rank of 79.7--128.1 against RankExt's final
cumulative delta's 44.9--66.4 at the same sampled layers. Arithmetic
averaging of five independently-trained deltas therefore neither
destructively cancels toward the single-specialist rank ($\sim$80) nor
purely stacks toward a naive ceiling ($5\times80=400$); it lands measurably
above a single specialist's rank, in the moderate-magnitude tail of the
spectrum specifically, not concentrated in the dominant directions (stable
rank is comparable between the two families). This measurement is real and
directly relevant, but it is not, on its own, an explanation of the
4.45-percentage-point gap. Two independent controlled capacity experiments
--- a rank-32-SimpleAvg-vs-cumulative-rank-160-RankExt comparison and a
same-seed rank-80-vs-160 doubling of RankExt's own cumulative rank at fixed
SimpleAvg configuration --- both found that increasing RankExt's capacity
within the tested range closed well under 5\% of the corresponding accuracy
gap for the non-KD variants. The results are consistent with a real,
measured effective-capacity edge for SimpleAvg's merge mechanism that is
nonetheless a minor, not dominant, contributor to the observed accuracy
gap; the larger share of the gap is better explained by the mechanisms
discussed in Sections~\ref{sec:ch6-rankext-struggles}--\ref{sec:ch6-kd-protect}
below.

One further, disclosed asymmetry is worth naming without overclaiming it.
SimpleAvg's classifier head trains at a learning-rate multiplier of
$\times10$ relative to the LoRA parameters, while RankExt's (in the primary
benchmark's configuration) trains at $\times1$ --- a deliberate,
historically-documented choice (Section~\ref{sec:ch6-calibration}), not an
unexplained oversight, but one that has never been ablated back to parity
for either family. This remains an open confound on the SimpleAvg-strength
question specifically, not resolved by any evidence gathered for this
thesis.

\section{Why Knowledge Distillation Does Not Improve SimpleAvg}
\label{sec:ch6-kd-simpleavg}

Chapter~\ref{ch:results} reports that SimpleAvg + KD (75.21\%) is, at this
single seed, effectively tied with plain SimpleAvg (75.16\%, $+$0.05
percentage points), while redistributing accuracy from the first task
($-$9.50 points) toward later tasks ($+$2.44 points) and raising restricted
accuracy slightly ($+$0.58 points). This section interprets that
redistribution.

SimpleAvg's KD objective \parencite{hinton2015distilling}
(Chapter~\ref{ch:methodology}) distills only the already-introduced
(``old-seen'') classes' logits from a frozen teacher onto the student's
current-task training images, at temperature 2, with no replay of old-task
images and with current- and future-task logits excluded from the
distillation softmax entirely. This design preserves the \emph{relative}
structure among old classes as reflected on new-task images --- it
constrains how the student's old-class logits relate to one another when
evaluated on images the student is currently training on --- but it does
not directly constrain old-class predictions on old-task images themselves,
since no old-task image ever appears in this loss term, and it places no
constraint at all on the total probability mass old classes receive
relative to newly introduced classes, since the current task's own logits
are masked out of the KD softmax by construction. This is consistent with,
and offers a plausible reading of, the observed pattern: restricted
accuracy (a within-task measure, insulated from cross-task competition)
improves slightly, and later-task accuracy improves, while first-task,
all-seen-scale, open accuracy --- which depends on exactly the
old-versus-new competition this KD objective does not constrain --- does
not.

A local gradient diagnostic on the SimpleAvg+KD configuration (one training
step, two optimizer updates) provides mechanistic support, at the tier of a
diagnostic, not a benchmark result. Immediately before the first optimizer
update, the KD gradient's norm is a substantial fraction of the
cross-entropy gradient's own norm (25.48 vs.\ 33.72, a ratio of roughly
three-quarters); after one and two optimizer updates the ratio grows
slightly ($23.15/28.96 \approx 80\%$; $21.65/27.75 \approx 78\%$). The
cosine similarity between the CE and KD gradients is small and positive
before any update ($+0.053$) and becomes mildly negative after both
subsequent updates ($-0.032$, then $-0.036$). This is consistent with a
real, non-trivial optimization interaction between the two loss terms,
strong enough to materially influence the update direction, but the
gradient conflict itself is mild, not catastrophic --- a cosine near zero,
not sharply negative. Pure CE/KD gradient opposition is therefore not, on
this evidence, a sufficient explanation on its own for the
first-task-versus-later-task redistribution observed in
Chapter~\ref{ch:results}; it is one plausible contributing factor among the
scope-of-constraint argument above, and the two are consistent with each
other rather than competing explanations.

This chapter does not conclude that SimpleAvg's KD objective is broken, nor
that it should be discarded; the redistribution it produces is a real,
measured trade-off (retention on later tasks against retention on the first
task), not an unambiguous failure. It also does not conclude that post-hoc
calibration is the cause of KD's flat all-seen result --- that question is
examined on its own terms in Section~\ref{sec:ch6-calibration}.

\section{SimpleAvg FactorOrth: Optimization Strength versus Geometric Objective}
\label{sec:ch6-sa-factororth}

Although both are referred to as FactorOrth in the thesis-facing
terminology, their concrete regularisation operators differ by family:
SimpleAvg applies orthogonality to dense updates, while RankExt applies it
to normalised low-rank factors (Section~\ref{sec:ch3-factororth}). This
section concerns the SimpleAvg, dense-update formulation only; RankExt's
factor-space FactorOrth is discussed in Section~\ref{sec:ch6-factororth}.

Chapter~\ref{ch:results} reports that neither SimpleAvg + FactorOrth
(74.02\%, $-$1.14 points versus plain) nor SimpleAvg + FactorOrth + KD
(73.22\%, $-$1.94 points versus plain) improves on the plain baseline, and
that lowering SimpleAvg FactorOrth's coefficient from $\lambda=20$ (the primary benchmark) to
$\lambda=1$ (the refinement experiment, Section~\ref{sec:ch5-ablation}) does not recover
the pure SimpleAvg FactorOrth variant's accuracy (73.80\% at $\lambda=1$, still
$-$0.22 points relative to $\lambda=20$'s 74.02\%, and still below plain).
This section separates two distinct questions the SimpleAvg FactorOrth result raises:
whether $\lambda=20$ was too strong an optimization pressure, and whether
the dense-update objective itself is well-aligned with the interference it is
intended to control.

On optimization strength: an independently-computed loss-magnitude audit,
based on an earlier diagnostic run's per-arm loss logs, finds that at $\lambda=20$,
the dense-update penalty's weighted loss contribution (in that run
applied to both families) is approximately 29.1\% of the cross-entropy loss
for RankExt, but only approximately 0.04\% (alone) to
0.15\% (combined with KD) of cross-entropy for SimpleAvg --- roughly three
orders of magnitude smaller. This is a loss-magnitude ratio, not a
gradient-norm ratio, and this chapter does not conflate the two: a
separate, SimpleAvg-specific gradient-level probe measured
weighted-orthogonality-to-CE ratios directly at candidate coefficients
between 0.75 and 2.0 (not 20) and found them likewise small throughout
training (median ratios of 0.001--0.033, maximum observed ratio 0.236 at
the single largest tested coefficient and earliest measured update). No
diagnostic file located for this thesis reports a gradient-norm ratio for
the dense-update penalty at $\lambda=20$ specifically for either family; a figure of this
kind was sought and not found, and is not used in this chapter. Taken
together, the available evidence --- the loss-magnitude ratio from an
earlier diagnostic run and the refinement experiment's $\lambda=1$
ablation both moving SimpleAvg FactorOrth's pure variant
no closer to, or further from, plain SimpleAvg --- is consistent with
the dense-update penalty being numerically near-inert for SimpleAvg across the coefficient
range actually tested (1 through 20), which argues against ``$\lambda=20$
was simply too aggressive'' as the explanation for its lack of benefit: an
inert term does not become effective by being made smaller.

On the geometric objective itself: SimpleAvg FactorOrth, as implemented, penalizes
each specialist's current dense weight update against every previous
specialist's own dense update, individually and pairwise. An alternative
formulation --- orthogonality against a merged or accumulated historical
direction, or against a shared subspace spanned by all previous updates,
rather than against each individual prior update separately --- was not the
mechanism tested in either the primary benchmark or the refinement experiment. This chapter does not
claim the individual-pairwise formulation is geometrically wrong; it states
it as a plausible limitation, distinct from and additional to the
coefficient question above: even a well-tuned coefficient could leave a
geometrically-misaligned penalty unable to constrain the interference it is
meant to target. Reducing $\lambda$ addresses only the first question (was
the pressure too strong) and, on the evidence in Chapter~\ref{ch:results}'s
ablation, does not resolve the second (is the objective itself targeting
the right geometric quantity) --- the two are logically separable, and
the refinement experiment's result speaks only to the first.

\section{Why Plain RankExt Struggles}
\label{sec:ch6-rankext-struggles}

Plain RankExt reaches 34.99\% all-seen accuracy against a restricted
accuracy of 91.03\% --- a 56.04-percentage-point gap
(Chapter~\ref{ch:results}, Section~\ref{sec:ch5-stability}), the largest of
any method in the benchmark. RankExt's architecture differs from
SimpleAvg's in exactly the properties Section~\ref{sec:ch6-simpleavg-strong}
identified as plausible contributors to SimpleAvg's strength: it is a
single, persistent adapter, sequentially trained, with previously-trained
rank blocks remaining active (not retrained, but live) in every later
forward pass, and a single classifier head shared across all five tasks
from the start. This chapter reports this as the structural difference
motivating RankExt's distinct failure mode, not as itself proving what that
failure mode is mechanistically.

The Chapter~\ref{ch:results} evidence directly rules out one candidate
explanation: insufficient adaptation capacity.
Section~\ref{sec:ch6-simpleavg-strong} already reports that two independent
controlled capacity experiments --- doubling RankExt's cumulative rank at
fixed protocol, and a rank32-SimpleAvg-vs-rank160-RankExt comparison in
which RankExt has substantially more nominal capacity than SimpleAvg ---
both find that increased capacity closes well under 5\% of the
corresponding gap for RankExt's non-KD variants, and in the plain case can
leave accuracy essentially unchanged or slightly worse. Within the tested
capacity range, capacity alone does not explain plain RankExt's weakness.
The open/restricted divergence itself
(Section~\ref{sec:ch6-open-restricted}) is more directly diagnostic:
restricted accuracy --- a within-task measure, insulated from competition
against other tasks' logits --- remains high (91.03\%) even as open,
all-seen accuracy collapses, which is inconsistent with a simple loss of
discriminative features and more consistent with a cross-task, output-space
competition effect, examined further below and in
Section~\ref{sec:ch6-open-restricted}. A companion audit of candidate
mechanisms found one further mechanism directly measured, on a related
variant set (not identical to the primary benchmark's own arms):
representation/feature drift, via a
cosine-alignment diagnostic comparing an old task's features to its own
classifier row versus the most recently trained row, strongly implicates
feature drift for non-KD RankExt configurations specifically. The primary benchmark's
non-KD arms (plain and FactorOrth-only) already include a training-time
mitigation for this --- a pretrained-feature anchor loss --- adopted in
response to that earlier diagnosis, meaning their reported
34.99\%/41.00\% results are already post-mitigation, not pre-fix numbers.

\section{The Role of KD+Protect in RankExt's Recovery}
\label{sec:ch6-kd-protect}

RankExt + KD + Protect reaches 65.01\% ($+$30.02 percentage points over
plain), the majority of the total recovery achieved by the combined
configuration (70.76\%). Following Chapter~\ref{ch:results}'s own
convention, this chapter attributes this recovery to the KD+Protect
configuration as a whole: full-scope, 100-way knowledge distillation and
the classifier-row feature-protection term are both active in every
RankExt arm that includes either, and this thesis's evidence does not
isolate their individual contributions from one another. This chapter does
not state that knowledge distillation alone causes this recovery.

Two properties of this configuration are consistent with why it addresses
RankExt's failure mode more directly than FactorOrth alone
(Section~\ref{sec:ch6-factororth}). Full-scope KD's teacher softmax spans
every class introduced so far, including the classes the student is
currently learning, so the distillation loss directly constrains the
\emph{relative scale} between old and newly introduced classes' logits ---
precisely the axis Section~\ref{sec:ch6-rankext-struggles}'s open/restricted
evidence implicates as the dominant failure mode, and precisely the
constraint SimpleAvg's differently-scoped, old-seen-only KD objective
(Section~\ref{sec:ch6-kd-simpleavg}) does not provide. The projected
feature-protection term penalizes the student's feature displacement, for
old classes, against a subspace derived from those classes' own classifier
rows --- a targeted retention constraint distinct in mechanism from KD's
output-space constraint. Both readings are offered as plausible
mechanistic accounts consistent with the measured result, not as
individually isolated causal findings; no ablation in this thesis separates
KD's contribution from Protect's.

\section{RankExt FactorOrth and KD+Protect}
\label{sec:ch6-factororth}

FactorOrth alone recovers 41.00\% from plain RankExt's 34.99\% ($+$6.01
percentage points) --- a real but much smaller recovery than KD+Protect's
$+$30.02 points. Combined with KD+Protect, the full recipe reaches 70.76\%,
a further $+$5.75 points above KD+Protect alone
(Chapter~\ref{ch:results}, Section~\ref{sec:ch5-rankext}). One plausible
interaction is that FactorOrth's weight-geometry regularization ---
penalizing the current rank block's factor-space cosine similarity against
previous blocks --- controls a different aspect of the adapter's structure
than KD+Protect's output-space and feature-space constraints, and that its
benefit is correspondingly larger when paired with a mechanism that already
addresses old-task retention directly. This is stated as an interaction
consistent with the measured pattern, not as a statistically established
synergy: this thesis ran no seed replication of any RankExt arm, and the
ablation evidence in Section~\ref{sec:ch6-calibration} shows this
interaction's sign is not stable across configurations (the
head-learning-rate refinement reverses it, Section~\ref{sec:ch6-calibration}).
FactorOrth alone, without KD+Protect, is not sufficient to close a
meaningful share of RankExt's gap to SimpleAvg --- its own-arm result
(41.00\%) remains far below SimpleAvg's weakest variant (73.22\%).

\section{Stability--Plasticity Trade-off}
\label{sec:ch6-stability-plasticity}

Chapter~\ref{ch:results}'s first-task, later-tasks-mean, BWT, and forgetting
values, its $R_{i,j}$ trajectories, and the head-learning-rate refinement's
combined-arm result together describe a consistent stability--plasticity
pattern \parencite{kirkpatrick2017overcoming,delange2022continual}:
mechanisms that improve retention of earlier tasks can reduce a model's
accuracy on tasks introduced later, and the balance is sensitive to exactly
how strongly a retention mechanism is applied, not only to whether one is
present. SimpleAvg's KD-bearing variants trade first-task accuracy for
later-task accuracy at roughly unchanged all-seen accuracy
(Section~\ref{sec:ch6-kd-simpleavg}). RankExt's auxiliary configurations
raise both first-task and later-task accuracy relative to plain, in an
order that matches their overall all-seen ranking --- auxiliary mechanisms
here improve both axes simultaneously, rather than trading one for the
other, which is a materially different pattern from SimpleAvg's.

The clearest single instance of this trade-off in this thesis is the
RankExt head-learning-rate refinement (the refinement experiment,
Section~\ref{sec:ch5-ablation}): raising the classifier-head learning-rate
multiplier from $\times1$ to $\times2$ improves three of the four RankExt
configurations, but the combined configuration --- the strongest result in
the primary benchmark --- regresses, from 70.76\% to 67.92\% ($-$2.84
percentage points). This regression is not a uniform degradation:
first-task accuracy improves ($60.70\% \to 67.75\%$, $+$7.05 points), BWT
improves ($-0.1584 \to -0.0532$), and forgetting improves ($0.2474 \to
0.1289$), while later-tasks-mean accuracy falls ($73.28\% \to 67.96\%$).
This is consistent with, and directly illustrates, the stability--plasticity
trade-off this section describes: at $\times2$, the combined configuration
becomes more stable on earlier tasks and less plastic on later ones, at
this single seed, and the net all-seen effect of that shift is negative for
this specific, most heavily retention-constrained configuration, even
though the same directional shift is net-positive for the three
configurations with fewer or weaker retention constraints. This chapter
interprets this as evidence that the optimal degree of classifier-head
plasticity for RankExt is not independent of which retention mechanisms
are simultaneously active in a given configuration --- raising plasticity
helps a lightly-constrained arm and can hurt a heavily-constrained one ---
rather than as evidence that either learning-rate setting is globally
preferable.

\section{The Local Head-Learning-Rate Probe and the Full Production Run}
\label{sec:ch6-lr-probe}

This section corrects an important point of scope that is easy to state
imprecisely. The local decision probe that informed the refinement experiment's choice of
a $\times2$ classifier-head learning-rate multiplier for RankExt tested the
\textbf{same RankExt method} used in the primary benchmark's flagship arm
--- the full FactorOrth + KD + Protect combined configuration --- not a
simpler or differently-composed RankExt variant. The final-compare script
that extended the sweep to LR$=1.0$ and LR$=2.0$ is confirmed, by direct
comparison, to be a byte-for-byte copy of the decision-probe script,
differing only in the swept learning-rate values, their assertions, and the
run-name prefix; both scripts target this one method exclusively. What
differs between the probe and the primary/ablation production runs is
\textbf{scale}, not \textbf{method}: the probe trained on 2 tasks (not 5),
2 epochs per task (not 9), and 10--15 images per class (not the full
dataset). Within that reduced scale, the probe found validation
cross-entropy and every measured accuracy metric improving monotonically as
the learning-rate multiplier increased across $\{0.1, 0.2, 0.5\}$ and then
across $\{0.5, 1.0, 2.0\}$, consistently across two seeds, with no
exception on any headline metric.

The full 5-task, 9-epoch, full-dataset production run (the refinement experiment) does not
reproduce this finding for the one configuration the probe actually tested:
the combined arm's accuracy falls under $\times2$
(Section~\ref{sec:ch6-stability-plasticity}). Three of the four RankExt
configurations --- none of which were the probe's own tested configuration
--- do improve under $\times2$, consistent with the probe's directional
finding extended by analogy, not by direct test, to those other
configurations. This is best read as a genuine methodological lesson about
probe scale, not about probe scope: a small-scale decision probe, run on
the exact production method it was meant to inform, correctly identified
the direction of the effect within the range it tested, but did not --- and
at that scale could not have been expected to --- predict that the same
method's response to a longer, full-length continual sequence, with five
tasks' worth of accumulated retention constraints active simultaneously
rather than two, would reverse sign for that specific,
already-most-heavily-constrained configuration. This chapter treats the
probe's own within-scale conclusion as sound and treats the discrepancy as
evidence that short, few-task decision probes may not extrapolate reliably
to longer continual sequences for a configuration that is already close to
its own stability--plasticity limit, rather than as evidence the probe
itself was flawed in design or target.

\section{Open versus Restricted Performance and Cross-Task Competition}
\label{sec:ch6-open-restricted}

The open/restricted divergence is the most consistent pattern across every
RankExt configuration in Chapter~\ref{ch:results}
(Section~\ref{sec:ch5-stability}): restricted accuracy stays in a narrow
91.03\%--94.78\% band across all eight methods in the benchmark, while
all-seen accuracy for RankExt alone spans 34.99\% to 70.76\%. Plain
RankExt's own gap (56.04 percentage points) is the largest observed. This
indicates, as Chapter~\ref{ch:results} states factually, that within-task
class discrimination remains substantially stronger than task-agnostic
discrimination for the affected configurations --- restricted accuracy, by
construction, measures whether the correct class is still favored once
competition from other tasks' classes is removed, and that measurement
stays high even where the full, unrestricted competition collapses the
same model's accuracy.

Multiple, not mutually exclusive, mechanisms are consistent with this
pattern, and this chapter does not conclude the gap is purely a calibration
artifact, nor that it demonstrates zero forgetting. Cross-task logit
competition --- newer tasks' logits systematically out-competing older
tasks' logits in the shared argmax --- is consistent with the monotone,
step-by-step decline visible in the $R_{i,j}$ matrices
(Chapter~\ref{ch:results}, Section~\ref{sec:ch5-rankext}): each old task's
open accuracy erodes as more tasks are trained, while its own diagonal
(just-learned) accuracy and its restricted accuracy both stay high
throughout. Score-scale mismatch between task groups is directly, though
diagnostically, evidenced by the post-hoc calibration result discussed in
Section~\ref{sec:ch6-calibration}. Representation drift is separately
implicated, for non-KD RankExt configurations specifically, by the
cosine-alignment diagnostic cited in
Section~\ref{sec:ch6-rankext-struggles}. These mechanisms plausibly coexist
rather than compete: a model can simultaneously retain within-task
discriminative structure, suffer some genuine representation drift for
non-KD-protected old tasks, and have its surviving structure obscured at
the output-space level by inter-task score misalignment. This thesis's
evidence does not decompose the 56.04-percentage-point plain-RankExt gap,
or any other configuration's smaller gap, into precise shares attributable
to each mechanism.

\section{Calibration as Diagnostic Headroom}
\label{sec:ch6-calibration}

A separate, explicitly diagnostic line of evidence bears on the
score-scale-mismatch reading above. A post-hoc, validation-fitted
inter-task logit-scale calibration \parencite{zhao2020maintaining} --- one
positive multiplicative scalar per task group, four free parameters, fit by
minimizing validation cross-entropy with no backbone, LoRA, or classifier
retraining and no test-set tuning --- was evaluated against an earlier,
separately-trained RankExt combined-arm checkpoint (6 epochs, no Protect30
term, not the primary benchmark's own weights, which were unavailable
locally for this diagnostic). On that checkpoint: raw, uncalibrated open
accuracy was
34.67\%; the pipeline's own historical row-norm calibration alone reached
53.12\% (matching that run's own reported result); adding the
four-parameter scale calibration reached 70.74\%, a further $+$17.62
percentage points, while restricted accuracy remained exactly unchanged at
93.68\% before and after (proven exactly invariant by construction, not
merely observed to be similar) --- direct evidence that a substantial
share of that specific checkpoint's residual open/restricted gap is
attributable to a correctable inter-task logit-scale mismatch, not to
further loss of within-task discriminative structure. The fitted scales
(2.147, 1.224, 1.178, 0.932, with the most recent task fixed at 1.0 as
reference) decrease roughly monotonically from oldest to most recent task,
consistent with a recency-competition reading: older tasks' logits needed
the largest upward correction to compete fairly against newer tasks'.

This chapter uses the term \textbf{post-hoc, validation-fitted inter-task
logit-scale calibration} throughout, and avoids ``classifier-weight
correction'' or ``corrected average weights'': the mechanism never touches
stored classifier weights, LoRA factors, or the merged dense delta; it
rescales the already-computed logit vector at inference time only, via a
small, separate module distinct from the model's own parameters. It is also
worth stating precisely what ``task-aware'' means here: the four scale
parameters are fit using knowledge of the fixed 20-class task-group
boundaries, but no per-sample task identity is required or used at
inference --- every test image is scored by the same, already-fitted,
group-wise rescaled classifier, regardless of which task it belongs to. The
mechanism is therefore task-structure-aware at calibration time and
task-agnostic at per-example inference time.

This result is reported strictly as diagnostic headroom, not as main
benchmark evidence, and for three compounding reasons: it was measured on
that earlier checkpoint, not the primary benchmark's own; that earlier
run lacks the Protect30 term and uses a different epoch count from the
primary benchmark, so its own 53.12\% starting point is not directly
comparable to the primary benchmark's 70.76\%; and the coincidence that
70.74\% (this diagnostic, on a weaker
starting checkpoint) lands close to the primary benchmark's own combined-arm result
(70.76\%, achieved through a different mechanism --- full-scope KD,
Protect30, and the pipeline's own more sophisticated row-norm calibration)
is noted as a striking but uncontrolled observation, not evidence that the
two mechanisms are interchangeable or additive. A related run, the calibration-diagnostic experiment,
applies this style of calibration (though not identically parameterized) to
RankExt while simultaneously changing SimpleAvg's classifier-head learning
rate; because the two changes are confounded within that run, the calibration-diagnostic experiment is
used in this chapter only as a second, independent diagnostic data point,
never as benchmark evidence: relative to the primary benchmark, it shows large gains
for RankExt configurations without a prior bias-correction mechanism
($+$32.71 and $+$31.47 percentage points for plain and FactorOrth-only
respectively) and small or negative changes for configurations that
already include KD+Protect ($+$1.75 and $-$1.17 percentage points for
KD+Protect and the combined configuration respectively) --- a pattern
consistent with calibration and KD+Protect addressing overlapping headroom
in the same underlying inter-task competition problem, though the calibration-diagnostic experiment's
own confound (the simultaneous SimpleAvg regression) prevents this from
being stated as a clean, isolated finding.

\section{Parameter and Rank Efficiency}
\label{sec:ch6-efficiency}

Chapter~\ref{ch:results}'s parameter-efficiency accounting
(Section~\ref{sec:ch5-efficiency}) shows the two families reach the same
final adapter size (a rank-80-equivalent, 2{,}949{,}120 LoRA parameters)
through structurally different training and deployment paths, and this
chapter does not conclude either is more efficient in general --- the two
represent different trade-offs, separable into training cost, storage, and
deployment overhead, each considered in turn.

\textbf{Training cost.} SimpleAvg trains 2{,}949{,}120 new LoRA parameters
at every one of its five tasks (a fresh, full rank-80 specialist each
time), for a cumulative training-time volume of 14{,}745{,}600 parameters
across the sequence. RankExt trains only 589{,}824 new parameters at each
task (a single new rank-16 block, with all earlier blocks frozen and never
retrained), for a cumulative training-time volume of 2{,}949{,}120 ---
exactly RankExt's own final adapter size, because each of its five blocks
is trained exactly once. RankExt is therefore substantially cheaper in
training-time parameter volume at every stage of the sequence, a genuine
structural property of its incremental design, independent of the accuracy
results discussed above.

\textbf{Deployment form.} The two families do not merely differ in size at
deployment; they differ in kind. SimpleAvg's merge step averages its five
specialists' dense per-module weight updates and adds the averaged update
directly into the pretrained base model's own weight tensors --- no
separate, low-rank adapter object survives this step. If the averaged
update is fused directly into the base weights, as the training pipeline
does, the deployed model requires no additional adapter parameters beyond
its 76{,}900-parameter classifier; the underlying rank-80-equivalent
structure that produced that update is not present in the deployed model
as a distinct, removable component. If, instead, the dense update were
stored separately from an unmodified base checkpoint --- for
reproducibility, versioning, or to keep a shared base model usable by
multiple downstream adapters --- it would occupy 14{,}155{,}776 values,
larger than the 2{,}949{,}120-parameter low-rank form it was derived from,
because the averaging step itself operates in dense, not factored, weight
space. This chapter does not treat the dense-storage figure as SimpleAvg's
default or mandatory deployment representation; it is one possible
representation, quantified because the fused, in-place form (SimpleAvg's
actual default) has no separately-quantifiable ``extra parameter'' figure
at all. RankExt's cumulative adapter, in contrast, remains in low-rank,
additive form at every step through deployment: its 2{,}949{,}120
parameters sit alongside an unmodified base model, addable and removable as
a distinct object, never fused into the base weights.

RankExt's structural advantage is therefore progressive, low training-time
cost and a persistent, separable adapter structure --- not, on this
evidence, a smaller final adapter than SimpleAvg's, since the two reach the
same nominal rank-80-equivalent size. SimpleAvg's structural advantage is a
final deployment representation that, once fused, requires no separate
adapter bookkeeping at all, at the cost of five times the training-time
parameter volume to reach it and no accessible low-rank structure at any
intermediate point in the sequence (a property directly connected to
Section~\ref{sec:ch5-rankext}'s observation that SimpleAvg has no genuine
$R_{i,j}$ trajectory: without a persistent, evolving intermediate model,
there is nothing to evaluate at ``after training through task $i$,
$i<5$''). No wall-clock or FLOPs measurement was made for either family in
this thesis; the comparison above is a parameter-count and
deployment-structure comparison only, and this chapter does not extend it
to computational cost.

\section{Limitations}
\label{sec:ch6-limitations}

The most consequential limitation is that every result presented as
primary in this thesis --- the primary benchmark, the ablation/refinement
comparison, and every number in this chapter drawn from either --- is a
\textbf{single-seed (seed 42)}
production result. No repeated-seed evidence exists for the primary
benchmark; the 4.45-percentage-point SimpleAvg-versus-RankExt gap, the
RankExt combined arm's $-$2.84-point response to the head-learning-rate
refinement, and every mechanistic interpretation offered above rest on this
single run each. A canonical multi-seed replication plan for the primary
benchmark exists but had not been executed at the time of writing.

Several of this chapter's auxiliary mechanisms are bundled, not
individually isolated: KD and Protect30 are never tested apart from one
another in RankExt (Section~\ref{sec:ch6-kd-protect}); FactorOrth's
interaction with KD+Protect is observed, not statistically isolated
(Section~\ref{sec:ch6-factororth}); the SimpleAvg FactorOrth loss-magnitude finding
and the SimpleAvg FactorOrth gradient-level probe come from different runs at
different coefficients, not a single controlled sweep through $\lambda=20$
itself. SimpleAvg's calibration behavior (a small-sample diagnostic: plain
accuracy moving from 84\% to 83\% with calibration on versus off, and the
KD variant moving from 83\% to 75\%, both on a 100-image evaluation subset)
is a real but tiny-sample observation that has not been run at benchmark
scale, and may confound comparisons between SimpleAvg configurations if
calibration interacts differently with each. BWT is constructed differently
for the two families by necessity (Section~\ref{sec:ch5-protocol}) and the
two families' BWT values, though tabulated together for compactness, are
not structurally comparable. SimpleAvg has no standard, persistent
$R_{i,j}$ trajectory to evaluate, for the structural reason given in
Section~\ref{sec:ch6-efficiency}, not a logging gap. The local
head-learning-rate probes and gradient/effective-rank diagnostics cited
throughout this chapter are, by their own design and by this thesis's own
evidence hierarchy, diagnostic only --- informative about direction and
mechanism, never substituted for benchmark evidence. A second dataset,
needed to test whether any conclusion in this chapter is specific to
CIFAR-100's particular class structure, had not been evaluated at the time
of writing. Finally, this thesis's parameter-efficiency comparison
(Section~\ref{sec:ch6-efficiency}) covers parameter counts and deployment
structure only; no wall-clock training time or FLOPs measurement was made
for either family, and no claim in this chapter should be read as a
computational-cost comparison.

\section{Future Work}
\label{sec:ch6-future-work}

The evidence gathered in this thesis motivates several specific,
evidence-grounded directions, stated in priority order rather than as a
generic list.

\textbf{Reformulate SimpleAvg FactorOrth's geometric target.}
Section~\ref{sec:ch6-sa-factororth}'s distinction between optimization strength
and geometric alignment motivates testing an orthogonality penalty against
a merged or accumulated historical direction, or against a subspace
spanned by all previous updates jointly, rather than against each
individual prior specialist's update pairwise --- a test that would
isolate whether the dense-update formulation, not only its coefficient, limits
its usefulness.

\textbf{Disentangle KD from Protect30 in RankExt.} No arm in this thesis
tests full-scope KD without the projected feature-protection term, or
Protect30 without KD, for RankExt; Section~\ref{sec:ch6-kd-protect}'s
attribution to ``the KD+Protect configuration'' as a whole is the most
precise claim this evidence supports, and resolving it further requires
this specific ablation.

\textbf{Revisit KD's data distribution for SimpleAvg, including
replay-aware variants.} Section~\ref{sec:ch6-kd-simpleavg}'s
scope-of-constraint argument suggests that a KD objective exposed to
old-task images directly (rather than only old-class logits on new-task
images), or an explicitly replay-aware distillation design, would more
directly constrain the axis Chapter~\ref{ch:results}'s evidence identifies
as unaddressed by the current old-seen, no-replay formulation.

\textbf{Complete calibration-aware continual learning as an integrated, not
diagnostic-only, mechanism.} Section~\ref{sec:ch6-calibration}'s headroom
result, measured only on a non-primary checkpoint, is a strong-tier
diagnostic finding but a weak-tier claim about the primary benchmark
specifically; the natural next step is fitting and evaluating the same
calibration mechanism against the primary benchmark's own checkpoints directly, and,
if it holds, integrating it as part of the trained or
immediately-post-training pipeline rather than a separate diagnostic pass.

\textbf{Multi-seed replication of the primary benchmark and the
head-learning-rate refinement.} Section~\ref{sec:ch6-limitations}'s
single-seed limitation applies to every quantitative claim in this thesis;
the combined-arm regression under the head-learning-rate refinement
(Sections~\ref{sec:ch6-stability-plasticity}--\ref{sec:ch6-lr-probe}) is
the single most novel and most seed-fragile finding, and is the
highest-priority candidate for replication.

\textbf{A second dataset}, to test whether the SimpleAvg-versus-RankExt
structural account developed in this chapter is a property of the two
architectures generally or is in some way specific to CIFAR-100's class
structure --- ImageNet-100 is a natural candidate; no results for a second
dataset exist at the time of writing.

\textbf{Protocol depth.} Every result in this thesis uses the primary
$5\times20$ protocol. Whether RankExt's behaviour, and its position
relative to SimpleAvg, changes when the same 100 classes are partitioned
into fewer or more incremental steps, with the total class count and final
cumulative rank held fixed (Section~\ref{sec:protocol-depth}), is not
examined here and is a natural controlled extension of this work.

\textbf{A retention-schedule or plasticity-aware learning-rate policy for
RankExt}, motivated directly by
Section~\ref{sec:ch6-stability-plasticity}'s finding that a single fixed
classifier-head learning rate does not suit every RankExt configuration
equally --- a schedule sensitive to how many retention mechanisms are
simultaneously active, rather than a single global multiplier, is a
plausible, evidence-motivated next step.

\textbf{Feature-space preservation as a complementary retention mechanism
for RankExt}, distinct from both KD's output-space constraint and
FactorOrth's weight-geometry constraint, is a plausible additional lever
suggested by, but not tested in, this thesis's mechanistic diagnostics
(Section~\ref{sec:ch6-rankext-struggles}'s representation-drift finding for
non-KD arms).

\section{Final Conclusion}
\label{sec:ch6-conclusion}

Within the evaluated setting --- CIFAR-100, $5\times20$ class-incremental
protocol, seed 42 --- this thesis's central empirical finding is an
asymmetry between two structurally different approaches to continual LoRA
adaptation. SimpleAvg, an architecturally simple strategy of independently
training per-task specialists and merging them once at the end, is
unexpectedly strong: it reaches 75.16\% without any auxiliary retention
mechanism, and none of the auxiliary mechanisms tested in this thesis ---
knowledge distillation, an orthogonality penalty, or their combination ---
improves on that baseline in any way this thesis can distinguish from
single-seed noise. RankExt, a persistent, incrementally-growing adapter
architecture, requires explicit retention mechanisms to approach
SimpleAvg's result at all: plain RankExt reaches only 34.99\%, and the
combination of full-scope knowledge distillation, classifier-row feature
protection, and a weight-geometry orthogonality penalty recovers most,
though not all, of the resulting gap, reaching 70.76\% --- 4.45 percentage
points below SimpleAvg's best result.

The evidence in this chapter is consistent with this asymmetry being
substantially structural rather than a matter of tuning either family
further within the settings already tested. RankExt's single, persistent,
shared classifier head exposes it to a cross-task, output-space
competition effect that its restricted-accuracy evidence shows is not
primarily a loss of underlying discriminative features; SimpleAvg's
independent-specialist, classifier-stitching design largely avoids this
competition by construction, at the cost of substantially higher
training-time parameter volume and a final deployment form that discards
its low-rank structure entirely. Auxiliary constraints that directly
target RankExt's output-space competition --- full-scope KD and
classifier-row protection --- recover the large majority of its deficit;
auxiliary constraints layered onto SimpleAvg's already-effective
architecture do not find an equivalent defect to repair, and in
the case of SimpleAvg FactorOrth appear numerically inert at every coefficient tested.
Cross-task score alignment remains a major, only partially resolved,
residual issue for RankExt: post-hoc calibration evidence, diagnostic and
not yet confirmed on the primary benchmark's own checkpoints, suggests a
substantial further share of its remaining gap is a correctable
scale-mismatch problem rather than additional representation damage.

More generally, this thesis's parameter-efficiency and
stability--plasticity evidence together indicate that parameter-efficient
continual adaptation via LoRA involves a genuine trade-off, not a single
dominant strategy, between independent-specialist merging and persistent,
incremental growth: the former is inexpensive to deploy but expensive to
train and architecturally blind to any notion of sequential state; the
latter is inexpensive to train incrementally and retains an inspectable,
persistent structure, but is exposed to cross-task interference that must
be actively managed, and the correct degree of that management is itself
sensitive to which other mechanisms are simultaneously active, as the
head-learning-rate refinement's mixed result demonstrates directly. This
thesis does not identify a universal winner between the two strategies
beyond the settings evaluated here; it identifies, with evidence at
multiple tiers, \emph{why} the two strategies fare differently under these
settings, and what would be required --- multi-seed replication, a second
dataset, and several specific, targeted follow-up experiments named in
Section~\ref{sec:ch6-future-work} --- to extend these conclusions further.
